%&../template
\documentclass[../main.tex]{subfiles}

\begin{document}
    \clearpage

    \section{Linux UIO}
    \begin{quote}
        \it
        For many types of devices, creating a Linux kernel driver is overkill. All that is really needed is some way to handle an interrupt and provide access to the memory space of the device. [...] Hardware that is ideally suited for an UIO driver fulfills all the following:

        \begin{itemize}
            \item The device has memory that can be mapped. The device can be controlled completely by writing to this memory.
            \item The device usually generates interrupts.
            \item The device does not fit into one of the standard kernel subsystems.
        \end{itemize}
    \end{quote}

    This fits very well for our usecase, the Mini-UART of the Raspberry Pi 4b, in particular because the device \enquote{has shallow FIFOs, no DMA Support}, can be controlled via Memory-Mapped-IO, meaning a simple \texttt{mmap} will suffice and it generates interrupts.


    \subsection{UIO Howto}

    \enquote{Each UIO device is accessed through a device file (\texttt{/dev/uioX}) and several sysfs attribute files.} 
    
    \subsubsection{Accessing Memory}
    \enquote{Each UIO device can make one or more memory regions available for memory mapping. Each mapping has its own directory in sysfs, the first mapping appears as \texttt{/sys/class/uio/uioX/maps/map0/}}

    \enquote{Each \texttt{mapX/} directory contains four read-only files that show attributes of the memory:}
    \begin{itemize}
        \item \texttt{name}: A string identifier for this mapping.
        \item \texttt{addr}: The address of memory that can be mapped.
        \item \texttt{offset}: The size, in bytes, of the memory pointed to by addr.
        \item \texttt{size}: The offset, in bytes, that has to be added to the pointer returned by mmap() to get to the actual device memory.
    \end{itemize}

    \enquote{From userspace, the different mappings are distinguished by adjusting the offset parameter of the mmap() call. To map the memory of mapping N, you have to use N times the page size as your offset}

    \subsubsection{Handling Interrupts}
    \enquote{Interrupts are handled by reading from /dev/uioX. A blocking read() from /dev/uioX will return as soon as an interrupt occurs. [...] The integer value read from /dev/uioX represents the total interrupt count}. This number should always be incrementing by 1.

    \enquote{UIO also implements a write() function. It is normally not used [...]. A write to \texttt{/dev/uioX} will call the \texttt{irqcontrol()} function implemented by the driver.} Writing a \texttt{1} will enable interrupts and writing \texttt{0} will disable them.

    \subsubsection{Inserting the UIO Module}
    Inserting the UIO Kernel module is straightforward, however there are a few things that require attention. First, we have to modify the device tree (\texttt{arch/arm/boot/dts/broadcom/bcm2710-rpi-3-b.dts}), and add an additional entry in the top level scope and disable the UART:
    \begin{verbatim}
/ {
    soc {
        uio_uart: uio_uart@7e215040 {
            compatible = "uio_uart";
            reg = <0x7e215000 0x6c>;     // Mini UART registers (address + size)
            interrupts = <1 29>;         // IRQ 29 with flags=1
            interrupt-parent = <&intc>;  // Interrupt controller reference
            status = "okay";
        };
    };
};

&uart0 {
    ...
    status = "disabled";
};
    \end{verbatim}

    \newpage
    Next, in our Kernel Module, we will have to specify the \texttt{compatible} string:

    \begin{verbatim}
static int uio_uart_probe(struct platform_device *pdev) {
    printk(KERN_INFO "Hello World!");
    TODO: This is not quite enough for /sys/class/uio/uio0 to appear...
    return 0;
};

static struct of_device_id uio_uart_match[] = {
    { .compatible = "uio_uart" },
    { /* Sentinel */ },
};

static struct platform_driver uio_uart = {
    .probe = uio_uart_probe,
    .remove = uio_uart_remove,
    .driver = {
        .name = DRIVER_NAME,
        .pm = &uio_uart_dev_pm_ops,
        .of_match_table = of_match_ptr(uio_uart_match),
    },
};
    \end{verbatim}

    With this setup done, we can finally load the module with \texttt{sudo insmod uio\_uart.ko} and see the following output in \texttt{dmesg}

\begin{verbatim}
    ...
    [ 1918.569519] Hello World!
\end{verbatim}

    The \texttt{struct uioinfo} tells the framework the details of our driver and is passed through\\\texttt{devm\_uio\_register\_device}.

    For local state in the driver, the UIO framework provides the \texttt{void *priv} pointer in the \texttt{uio\_info} struct. We will simply allocate a \texttt{sizeof(struct uio\_uart\_privdata *)} and store it there. Upon entering a handler function, the pointer to \texttt{uio\_info} will be passed, which contains the \texttt{priv} pointer.


    % All init code should go to userspace
    %   "kernel part is just a thin wrapper to give userspace safe access to device memory and interrupts"


\end{document}